\documentclass[12pt]{article}
\usepackage{amsmath, amssymb}
\usepackage{geometry}
\usepackage{hyperref}
\usepackage{enumitem}
\usepackage{array}
\geometry{margin=1in}

%--------------------
\begin{document}

\begin{center}
\textbf{\Large Project 2 \\ How Reliable Is Linear Regression Inference?} \\
\end{center}


\section*{Overview}

Classical linear regression inference --- confidence intervals and hypothesis tests for slopes --- rests on a set of assumptions about the error distribution: normality, constant variance, and independence. In practice, real data rarely satisfy all of these cleanly. This project investigates those questions:
\begin{itemize}
    \item How badly do classical 95\% confidence intervals for regression slopes fail when assumptions are violated?
    \item How p-value can be used in the bootstrap when assumptions are violated?
    \item Does the bootstrap provide a reliable remedy?
\end{itemize}
You will answer these questions in two ways. In Part I, you will use simulation --- where the true parameter $\beta_1 = 2$ is known --- to directly measure how often each CI method works under a specific type of violation. In Part II, you will apply both methods to a real dataset and use your simulation results to assess how much trust to place in each approach. \\

\noindent\textbf{Each group will choose two of four simulation studies for Part I (Studies A--D).} \\

\noindent\textbf{Your group will submit a written report (8--12 pages, excluding code appendix) and a complete, reproducible R script.}


\section*{Suggested Outline for the Report}
\begin{itemize}
    \item Introduction
    \begin{enumerate}
        \item Motivation
        \item Which two simulation studies your group chose, and why those violations are interesting
        \item Overview of the two CI methods being compared (classical OLS vs.\ bootstrap)
        \item Overview of the two p-value methods being compared (classical $t$-test vs.\ bootstrap)
        \item Brief preview of findings
    \end{enumerate}

    \item Methods
    \begin{enumerate}
        \item Base data-generating model and simulation design
        \item Description of your two chosen studies and their scenarios
        \item Definition of empirical coverage probability, average CI width, and bias
        \item Definition of the bootstrap p-value for testing $H_0\colon \beta_1 = 0$ (null-centered percentile method)
        \item Description of the real dataset chosen for Part II
    \end{enumerate}

    \item Part I: Simulation Study Results
    \begin{enumerate}
        \item Empirical coverage probabilities (table: 4 scenarios $\times$ 2 sample sizes $\times$ 2 CI methods, for each study)
        \item Average CI widths
        \item Bias of $\hat{\beta}_1$ across scenarios
        \item Empirical rejection rates for $H_0\colon \beta_1 = 0$: classical $t$-test p-value vs.\ bootstrap p-value (table: same structure as coverage table)
        \item Histograms of the sampling distribution of $\hat{\beta}_1$ (one per scenario)
        \item Scenario-by-scenario interpretation: what drives the pattern you observe?
        \item Comparison across your two studies: which violation is more damaging? Does the bootstrap help differently?
    \end{enumerate}

    \item Part II: Real Data Analysis
    \begin{enumerate}
        \item Exploratory data analysis and motivation for chosen predictor
        \item OLS model output, classical 95\% CI, and classical p-value for the slope
        \item Bootstrap 95\% CI, bootstrap p-value for $H_0\colon \beta_1 = 0$, and bootstrap distribution plot
        \item Residual diagnostics (residual plot and Q-Q plot)
        \item Comparison of the two CIs and the two p-values: do they agree? Are assumptions likely violated?
        \item Connection to Part I: which of your studied scenarios does this most resemble, and what does that imply about how much to trust the classical CI and p-value?
    \end{enumerate}

    % \item Part III: Reflection
    % \begin{enumerate}
    %     \item When does the bootstrap help (or not help) for your chosen violation types?
    %     \item Does larger sample size solve the problems your studies identified?
    %     \item What would you recommend to a researcher who suspects these violations in their data?
    % \end{enumerate}

    \item Conclusion
    \begin{enumerate}
        \item Summary of key findings from simulation and real data
        \item Practical guidance specific to your chosen violation types
        \item Suggestions for future work
    \end{enumerate}
\end{itemize}


\section*{Part I: Simulation Study}

\subsection*{Base Setting}

\textbf{Base data-generating model:}
\[
Y_i = \beta_0 + \beta_1 X_i + \epsilon_i, \quad \beta_0 = 1,\quad \beta_1 = 2,\quad X_i \sim N(0,1)
\]
The only thing that varies across groups is the distribution of $\epsilon_i$.

\begin{itemize}
    \item \textbf{Sample sizes:} different $n$ for each scenario, e.g., $n = 30, 100, 200$
    \item \textbf{Replications:} $500$ or $1{,}000$ simulated datasets per scenario (bootstrap sample size $B = 1{,}000$)
    \item \textbf{CI methods:} Classical 95\% CI and bootstrap 95\% percentile CI
    \item \textbf{P-value methods:} Classical $t$-test p-value and bootstrap p-value for $H_0\colon \beta_1 = 0$, computed by centering the bootstrap distribution of $\hat{\beta}_1^*$ at zero and finding the proportion of resamples satisfying $|\hat{\beta}_1^* - \bar{\hat{\beta}}_1^*| \geq |\hat{\beta}_1|$ (two-sided)
\end{itemize}

\subsection*{For Each Cell, Report:}
\begin{itemize}
    \item Empirical coverage probability (nominal target: 0.95) for both CI methods
    \item Average CI width for both methods
    \item Bias: $\overline{\hat{\beta}}_1 - 2$
    \item Empirical rejection rate for $H_0\colon\beta_1 = 0$ using both the classical and bootstrap p-values 
    \item Histogram of the sampling distribution of $\hat{\beta}_1$
\end{itemize}

%--------------------------------------------------
\subsection*{Study A --- Non-normality of Errors}

\textbf{Violation:} The classical CI assumes normally distributed errors. This study investigates how skewness in the error distribution affects coverage, from mild to severe. All errors are mean-centered so $E(\epsilon_i) = 0$ and the model remains correctly specified.

%--------------------------------------------------
\subsection*{Study B --- Heteroskedasticity}

\textbf{Violation:} OLS assumes $\mathrm{Var}(\epsilon_i) = \sigma^2$ (constant variance). This study investigates what happens when variance grows with $|X_i|$.

%--------------------------------------------------
\subsection*{Study C --- Heavy-Tailed Errors}

\textbf{Violation:} Fat-tailed errors increase the frequency of extreme residuals and inflate the true variability of $\hat{\beta}_1$ beyond what the classical SE estimates. Tail weight is controlled by the degrees of freedom $\nu$ of a $t$-distribution; smaller $\nu$ means heavier tails.

%--------------------------------------------------
\subsection*{Study D --- Model Misspecification}

\textbf{Violation:} The fitted model is linear ($\hat{Y} = \hat{\beta}_0 + \hat{\beta}_1 X$), but the true relationship contains a quadratic component. Unlike Studies A--C, this is a \emph{model misspecification} violation rather than a distributional one. It introduces \emph{bias} in $\hat{\beta}_1$ in addition to coverage problems, because the linear slope absorbs part of the curvature.


\section*{Part II: Real Data Analysis}

Choose \textbf{one} dataset from the suggested sources below, or propose your own. Your dataset must have a continuous response variable and at least one continuous predictor with a plausible linear relationship. In your discussion, connect the assumption violations you detected in the data to the specific violation types your group investigated in Part I.

\subsection*{Required Steps}
\begin{enumerate}
    \item Perform brief EDA and motivate your choice of primary predictor.
    \item Fit OLS and report the classical 95\% CI and classical p-value for the slope.
    \item Compute the bootstrap 95\% CI ($B = 1{,}000$, percentile method).
    \item Compute the bootstrap p-value for $H_0\colon \beta_1 = 0$ ($B = 1{,}000$, null-centered percentile method). Report and interpret it alongside the classical p-value.
    \item Plot the bootstrap distribution of $\hat{\beta}_1$ and describe its shape.
    \item Check model assumptions using a residual plot and a Q-Q plot.
    \item Discuss: do the two CIs agree? Do the classical and bootstrap p-values agree? Are assumptions likely violated? Which of your Part I scenarios does this most resemble?
    \item Based on your simulation results, how much do you trust the classical CI and p-value here?
\end{enumerate}


\section*{Project Requirements}
\begin{itemize}
    \item Source of data (Part II) must be cited
    \item Complete R code must be included as an appendix (\texttt{.R} or \texttt{.Rmd})
    \item \texttt{set.seed()} must be used for all simulations to ensure reproducibility
    \item References should be provided at the end of the report
\end{itemize}


\section*{Possible Datasets for Part II}
\begin{itemize}
    \item \textbf{Boston Housing} --- median home value vs.\ socioeconomic predictors (\texttt{MASS} R package)
    \item \textbf{Auto MPG} --- fuel efficiency vs.\ vehicle characteristics (\href{https://archive.ics.uci.edu/}{UCI ML Repository})
    \item \textbf{Gapminder} --- life expectancy vs.\ GDP per capita (\texttt{gapminder} R package)
    \item \href{https://www.openintro.org/data/}{OpenIntro Datasets}
    \item \href{https://www.kaggle.com/datasets}{Kaggle Datasets}
    \item \href{https://data.gov/}{US Government Open Data}
\end{itemize}


\section*{Deliverables}
\begin{itemize}
    \item Deliverable 1: Project Outline due on \textbf{Tuesday Mar 17 at 10:00 PM}
    \item Deliverable 2: Progress Report (Draft slides + R code + dataset) due on \textbf{Monday March 30 at 10:00 PM}
    \item Deliverable 3: Final Written Report (PDF) + R script + dataset due on \textbf{Tuesday Apr 7 at 10:00 PM}
\end{itemize}


\section*{Important Dates (Tentative)}
\begin{itemize}
    \item \textbf{Wednesday March 4 Lab:} Project 2 logistics --- study selection, dataset selection, and simulation design
    \item \textbf{Tuesday Mar 17 at 10:00 PM:} Deliverable 1 --- Project Outline due
    \item \textbf{Wednesday March 18 Lab:} Project 2 work lab --- Part I simulation
    \item \textbf{Wednesday March 25 Lab:} Project 2 work lab --- Part II real data analysis
    \item \textbf{Monday March 30 at 10:00 PM:} Deliverable 2 --- Progress Report due
    \item \textbf{Wednesday April 1 Lab:} Project 2 work lab --- reflection and report finalisation
    \item \textbf{Tuesday Apr 7 at 10:00 PM:} Deliverable 3 --- Final Report + R script due
    \item \textbf{Wednesday Apr 29:} Resubmission deadline (for groups receiving R or N)
\end{itemize}


\end{document}
